\section{La ventana de aplicaciones de 2025: por qué Lovart se dio a conocer más allá de su nicho en este momento}

Que Lovart se diera a conocer más allá de su nicho no es un hecho aislado: ocurrió después de que en 2025 se abriera la ventana para las aplicaciones de Agent. Los modelos base ya son lo bastante potentes y los usuarios empiezan a entender el valor de los Agent, pero los Agent de propósito general todavía no cubren todos los flujos de trabajo profesionales. Esta ventana deja una oportunidad a las aplicaciones verticales: si logran ofrecer en un ámbito profesional una experiencia claramente superior a la de las herramientas de propósito general, pueden obtener con rapidez atención a escala mundial.

\begin{knowledgebox}{Las tres condiciones de la ventana de aplicaciones}
\begin{tabularx}{\textwidth}{p{0.20\textwidth}p{0.34\textwidth}X}
\toprule
Condición & Explicación & Correspondencia en Lovart \\
\midrule
Capacidad suficiente del modelo & El modelo base puede entender la necesidad y generar contenido de alta calidad & La generación de diseños y la modificación en varias rondas ya son utilizables. \\
Los puntos de entrada generales no lo abarcan todo & ChatGPT/Manus todavía no pueden profundizar en todos los ámbitos verticales & Sigue habiendo un espacio libre en el flujo de trabajo de los diseñadores. \\
Percepción madura del usuario & Los usuarios están dispuestos a intentar delegar tareas profesionales a la IA & Diseñadores y creadores empiezan a buscar un entorno de trabajo con IA. \\
\bottomrule
\end{tabularx}
\end{knowledgebox}

\begin{warningbox}{La ventana de oportunidad también es una cuenta regresiva}
Cuando un Agent vertical despega, los Agent de propósito general, las empresas de modelos y los competidores aprenden de él. Al mismo tiempo que la ventana se abre, también se va cerrando. Lovart debe consolidar su flujo de trabajo y su lugar en la percepción de los usuarios antes de que los demás lo imiten.
\end{warningbox}

\subsection{Resumen de la sección}

Que Lovart se diera a conocer más allá de su nicho se debe a la suma de la capacidad de los modelos, la percepción madura de los usuarios y un espacio libre en un ámbito vertical. Su oportunidad no es permanente: debe convertir cuanto antes la ventaja de la ventana en una barrera basada en el flujo de trabajo.
